\section{Physical correspondences, affine fibres, and torsion}
\label{sec:states-filtration}

The structural part of the proof has three ingredients: compatible interval configurations with additive observations, finite labelled correspondences composed by pullback over the literal compatibility domain, and affine fibres in the prime-power filtration of $\Q/\Z$.  The constructions below are used only where a universal property controls a later transport or pasting step.

Let $\mathcal C$ be the set of finite subsets $S\subset\N_{>0}$ whose connected components in the path graph have cardinality $2$ or $3$.  Put
\begin{equation}\label{eq:observations}
W(S)=\sum_{n\in S}\frac1n,
\qquad
\nu(S)=|\pi_0(S)|,
\qquad
\overline W(S)=[W(S)]\in A:=\Q/\Z.
\end{equation}
Let
\[
M=A\times\N,
\qquad
0_M=(0,0),
\qquad
\Omega=(\overline W,\nu):\mathcal C\longrightarrow M.
\]
For $n\ge1$, let
\[
D_n\subseteq\mathcal C^n
\]
be the joint-compatibility domain: $(S_1,\ldots,S_n)\in D_n$ when the supports are pairwise disjoint and their literal union is again an element of $\mathcal C$.  Write $\mu_n:D_n\to\mathcal C$ for literal union.  Hence
\begin{equation}\label{eq:additive-observations}
W\!\left(\bigcup_iS_i\right)=\sum_iW(S_i),
\quad
\nu\!\left(\bigcup_iS_i\right)=\sum_i\nu(S_i),
\quad
\overline W\!\left(\bigcup_iS_i\right)=\sum_i\overline W(S_i).
\end{equation}
Equivalently, for binary compatible unions the square
\begin{equation}\label{eq:observation-square}
\begin{tikzcd}
D_2 \arrow[r,"{\mu_2}"] \arrow[d,"{(\Omega\pi_1,\Omega\pi_2)}"']
  & \mathcal C \arrow[d,"{\Omega}"]\\
M^2 \arrow[r,"{+}"'] & M
\end{tikzcd}
\end{equation}
commutes.

\begin{lemma}[Compatibility substitution law]\label{lem:compatibility-substitution}
Every subtuple of a point of $D_n$ is jointly compatible.  More generally, if the indices are partitioned into blocks, taking the union within each block gives a compatible tuple of block unions, and its final union is $\mu_n$.  Conversely, every iterated binary-compatible parenthesization determines a point of $D_n$.  Thus all parenthesizations have the common domain $D_n$ and the common value $\mu_n$.
\end{lemma}

\begin{proof}
Suppose $(S_1,\ldots,S_n)\in D_n$.  If a component of $S_i$ were adjacent to a component of $S_j$ with $i\ne j$, those two components would lie in one component of the total union and that component would have cardinality at least $4$, contrary to the definition of $\mathcal C$.  Components belonging to different labels are therefore separated by at least one unused integer.  Every subunion and every block union is consequently again an element of $\mathcal C$, and substitution does not change the literal total union.

Conversely, an iterated binary-compatible product records pairwise disjoint labels whose final literal union belongs to $\mathcal C$; these are exactly the defining conditions for membership in $D_n$.
\end{proof}

Thus $\mathcal C$, with unit $\varnothing$, domains $D_n$ and partial products $\mu_n$, is a partial commutative monoid, and $\Omega$ is additive wherever the partial product is defined.

\begin{remark}[Abstract scope]\label{rem:abstract-labelled-relations}
Everything in the next two subsections is formal for a partial commutative monoid $\mathscr C$ equipped with an additive map $\Omega:\mathscr C\to M$ to a commutative monoid and a nonnegative additive mass $W$.  The interval model enters only through the verification of the partial-sum axioms and, later, through the arithmetic estimates and the incompatibility graph.
\end{remark}

\subsection{Witnesses and regular images}

All pullbacks and images below are taken in $\mathbf{Set}$.  A regular epimorphism is therefore a surjection; its stability under pullback and composition is the mechanism that transports existence.  We retain a finite witness object $E$ and its target map $E\twoheadrightarrow Y$ until compatibility has been imposed.  Passing to an image records which targets occur but may discard the information needed to compose their witnesses.

\subsection{Labelled correspondences and compatible pasting}

Let $X,Y\subseteq M$ be finite.  A \emph{labelled finite correspondence}
\[
R:X\rightsquigarrow Y
\]
consists of a finite edge object $E_R$ and maps
\[
s_R:E_R\to X,
\qquad
\ell_R:E_R\to\mathcal C,
\qquad
t_R:E_R\to Y
\]
satisfying the observable conservation law
\begin{equation}\label{eq:edge-conservation}
t_R(e)=s_R(e)+\Omega(\ell_R(e))
\qquad(e\in E_R).
\end{equation}
Equivalently, the square
\begin{equation}\label{eq:edge-conservation-square}
\begin{tikzcd}
E_R \arrow[r,"{(s_R,\ell_R)}"] \arrow[d,"{t_R}"']
  & X\times\mathcal C \arrow[d,"{(x,L)\mapsto x+\Omega(L)}"]\\
Y \arrow[r] & M
\end{tikzcd}
\end{equation}
commutes, where the bottom arrow is inclusion.  The correspondence is \emph{covering} when $t_R:E_R\twoheadrightarrow Y$ is a regular epimorphism.

For correspondences $R:X\rightsquigarrow Y$ and $S:Y\rightsquigarrow Z$, first form the path pullback
\begin{equation}\label{eq:path-pullback}
\begin{tikzcd}
P(S,R) \arrow[r,"{\mathrm{pr}_S}"] \arrow[d,"{\mathrm{pr}_R}"']
  & E_S \arrow[d,"{s_S}"]\\
E_R \arrow[r,"{t_R}"'] & Y.
\end{tikzcd}
\end{equation}
The pair of label maps gives
\[
(\ell_R\mathrm{pr}_R,\ell_S\mathrm{pr}_S):P(S,R)\longrightarrow\mathcal C^2.
\]
The compatible path object is the pullback
\begin{equation}\label{eq:compatibility-pullback}
\begin{tikzcd}
P^{\mathrm{comp}}(S,R) \arrow[r] \arrow[d]
  & D_2 \arrow[d]\\
P(S,R) \arrow[r,"{(\ell_R\mathrm{pr}_R,\ell_S\mathrm{pr}_S)}"']
  & \mathcal C^2.
\end{tikzcd}
\end{equation}
Define $S\odot R$ to have edge object $P^{\mathrm{comp}}(S,R)$, source $s_R\mathrm{pr}_R$, target $t_S\mathrm{pr}_S$, and label obtained from the upper map in \eqref{eq:compatibility-pullback} followed by union $D_2\to\mathcal C$.  The conservation square for $S\odot R$ is the pasted composite of \eqref{eq:observation-square}, \eqref{eq:path-pullback}, and the two conservation squares of $R$ and $S$.

The identity correspondence on $X$ has edge object $X$ and empty label.  A $2$-cell is a map of edge objects preserving source, target and label; vertical composition is composition of these edge maps.  If $\alpha:R\Rightarrow R'$ and $\beta:S\Rightarrow S'$, the induced map $P(S,R)\to P(S',R')$ restricts, by the universal property of the compatibility pullback, to $P^{\mathrm{comp}}(S,R)\to P^{\mathrm{comp}}(S',R')$; this is horizontal composition.  This is the only higher structure required here: it records changes of witnesses without changing their observations or their physical labels.

\begin{lemma}\label{lem:labelled-relation-category}
Compatible pasting makes the finite state sets and labelled correspondences into a bicategory $\mathbf{Corr}_{\mathcal C}^{\Omega}$.  More precisely, both parenthesizations of any finite composite are canonically represented by the same compatible path object
\[
P_n^{\mathrm{comp}}
=
P_n\times_{\mathcal C^n}D_n,
\]
where $P_n$ is the iterated fibre product of the edge objects over the intermediate state sets.  We henceforth suppress these canonical pullback associators.
\end{lemma}

\begin{proof}
The iterated path object $P_n$ is characterized by the universal property of the fibre product, independently of parenthesization.  Pulling it back along $D_n\hookrightarrow\mathcal C^n$ imposes joint compatibility in one step.  Lemma~\ref{lem:compatibility-substitution} shows that every intermediate block union required by any parenthesization is defined on this same object and has total label $\mu_n$.  The identity statement follows from the empty label.  The canonical pullback isomorphisms satisfy the usual coherence because they are uniquely determined by their projections.
\end{proof}

The underlying unlabelled relation is the regular image
\begin{equation}\label{eq:underlying-relation}
U(R)=\operatorname{im}(s_R,t_R)\subseteq X\times Y.
\end{equation}
A $2$-cell $R\Rightarrow R'$ induces the inclusion $U(R)\subseteq U(R')$, which is the action of $U$ on $2$-cells.
For $E_R\ne\varnothing$ put
\[
\omega(R)=\max\{W(\ell_R(e)):e\in E_R\},
\]
and set $\omega(R)=0$ when $E_R=\varnothing$.

Since \eqref{eq:compatibility-pullback} is a subobject of the unrestricted path pullback, every composable pair has a canonical comparison
\begin{equation}\label{eq:oplax-inclusion}
U(S\odot R)\subseteq U(S)\circ U(R),
\end{equation}
The bicategory $\mathbf{Rel}$ has sets as objects, relations as arrows and inclusions as $2$-cells.  Thus $U$ is a normal oplax homomorphism to $\mathbf{Rel}$, with the comparison direction displayed in \eqref{eq:oplax-inclusion}.  The possible failure of equality has a concrete meaning: two successive edges may reuse the same nonempty interval label.  Their endpoint observations compose, but their labels do not.

For a finite composable string $R_1,\ldots,R_n$, let
\[
\pi_{0,n}:P_n\longrightarrow X_0\times X_n
\]
be the endpoint map and let $\ell:P_n\to\mathcal C^n$ be the label map.  Then
\begin{equation}\label{eq:strictness-endpoint-images}
U(R_n\odot\cdots\odot R_1)
=\operatorname{im}\bigl(\pi_{0,n}|_{P_n^{\mathrm{comp}}}\bigr),
\qquad
U(R_n)\circ\cdots\circ U(R_1)
=\operatorname{im}(\pi_{0,n}).
\end{equation}
Hence the comparison is an identity precisely when
\begin{equation}\label{eq:strictness-endpoint-criterion}
\operatorname{im}\bigl(\pi_{0,n}|_{P_n^{\mathrm{comp}}}\bigr)
=
\operatorname{im}(\pi_{0,n}).
\end{equation}
The information loss is displayed by the commuting image square
\begin{equation}\label{eq:strictness-information-square}
\begin{tikzcd}
P_n^{\mathrm{comp}} \arrow[r] \arrow[d,"{\pi_{0,n}}"']
  & P_n \arrow[d,"{\pi_{0,n}}"]\\
U(R_n\odot\cdots\odot R_1) \arrow[r]
  & U(R_n)\circ\cdots\circ U(R_1),
\end{tikzcd}
\end{equation}
where the lower objects denote the corresponding regular images.  A stronger witnesswise condition is that the path-label map factor through the compatibility domain:
\begin{equation}\label{eq:compatibility-factorization}
\begin{tikzcd}
P_n \arrow[r,"{\widetilde\ell}"] & D_n \arrow[r,"{\iota}"] & \mathcal C^n,
\end{tikzcd}
\qquad
\ell=\iota\circ\widetilde\ell.
\end{equation}
By the universal property of \eqref{eq:compatibility-pullback}, this factorization is equivalent to the canonical map
\[
P_n^{\mathrm{comp}}\longrightarrow P_n
\]
being an isomorphism.  It therefore implies
\begin{equation}\label{eq:rel-strictification}
U(R_n\odot\cdots\odot R_1)
=
U(R_n)\circ\cdots\circ U(R_1).
\end{equation}
Moreover, for every composable string,
\begin{equation}\label{eq:mass-subadditive}
\omega(R_n\odot\cdots\odot R_1)
\le\sum_{i=1}^n\omega(R_i),
\end{equation}
because the label of a compatible path is the literal union of its labels and reciprocal value is additive there.

\begin{theorem}[Compatible covering theorem]\label{thm:compatible-covering}
Let $R_i:X_{i-1}\rightsquigarrow X_i$ be a finite string of covering labelled correspondences.  If its path-label map factors through $D_n$, then the compatible composite is covering.  Its target map is obtained inductively from the local target maps by pullback and composition, and its label mass is at most $\sum_i\omega(R_i)$.
\end{theorem}

\begin{proof}
Let $P_i$ be the path object through $R_i$, and let $\tau_i:P_i\to X_i$ be its target map.  For $i\ge2$ the square
\[
\begin{tikzcd}
P_i \arrow[r,"{\mathrm{pr}_i}"] \arrow[d]
  & E_{R_i} \arrow[d,"{s_i}"]\\
P_{i-1} \arrow[r,"{\tau_{i-1}}"'] & X_{i-1}
\end{tikzcd}
\]
is a pullback.  Starting from $\tau_1=t_1$, pullback stability makes $\mathrm{pr}_i$ a regular epimorphism, and composition makes $\tau_i=t_i\mathrm{pr}_i$ a regular epimorphism.  The compatibility factorization identifies $P_n^{\mathrm{comp}}$ with $P_n$.  Under this identification its target is $\tau_n$, and its label is $\mu_n\widetilde\ell$.  The mass assertion is \eqref{eq:mass-subadditive}.
\end{proof}

The theorem separates two proof obligations.  Quotient calculations must provide covering local correspondences; the physical construction must make their path labels factor through $D_n$.  Once these obligations are met, the global existence statement follows by pasting pullbacks.  Neither obligation can be replaced by the other.

\subsection{Finite torsion and affine fibres in \texorpdfstring{$\Q/\Z$}{Q/Z}}

For $n\ge1$ put
\[
H_n=\{x\in A:nx=0\}.
\]
For a subgroup $H\le A$, let
\[
q_H:A\longrightarrow A/H
\]
be the quotient map, and for $\alpha\in A/H$ write
\[
\operatorname{Fib}_H(\alpha)=q_H^{-1}(\alpha)
\]
for the corresponding affine torsor fibre.  When $H$ is finite, this fibre is a finite coset of $H$.  In particular, every set $\operatorname{Fib}_H(\alpha)\times I$ used below is a finite state object.

\subsection{Affine extension as a pullback of regular epimorphisms}

Let $H\le K\le A$ be finite subgroups and write
\[
\pi_{H,K}:A/H\longrightarrow A/K
\]
for the quotient map.  The subgroup inclusion is encoded by the exact sequence
\begin{equation}\label{eq:affine-extension-exact}
\begin{tikzcd}[column sep=small]
0 \arrow[r] & K/H \arrow[r] & A/H \arrow[r,"{\pi_{H,K}}"] & A/K \arrow[r] & 0
\end{tikzcd}
\end{equation}
and by the pullback square
\begin{equation}\label{eq:affine-extension-pullback}
\begin{tikzcd}
K \arrow[r] \arrow[d,"{q_H|_K}"'] & A \arrow[d,"{q_H}"]\\
K/H \arrow[r] & A/H.
\end{tikzcd}
\end{equation}
For $\beta\in A/K$ put
\[
\operatorname{Fib}_{\pi_{H,K}}(\beta)
=\pi_{H,K}^{-1}(\beta)\subseteq A/H.
\]
Translation by $\alpha\in A/H$ gives a canonical torsor isomorphism
\begin{equation}\label{eq:fibre-translation-isomorphism}
\operatorname{Fib}_{\pi_{H,K}}(\beta)
\xrightarrow{\;u\mapsto\alpha+u\;}
\operatorname{Fib}_{\pi_{H,K}}\!\bigl(\pi_{H,K}(\alpha)+\beta\bigr).
\end{equation}

\begin{lemma}[Affine torsor extension principle]\label{lem:affine-torsor-extension}
Let $\alpha\in A/H$ and $\beta\in A/K$.  Let $d:Y\to A$ be a map such that $q_Kd$ is constantly $\beta$ and
\begin{equation}\label{eq:affine-increment-regular-epi}
q_Hd:Y\twoheadrightarrow\operatorname{Fib}_{\pi_{H,K}}(\beta)
\end{equation}
is a regular epimorphism.  Put
\[
F^- = \operatorname{Fib}_H(\alpha),
\qquad
F^+ = \operatorname{Fib}_K\!\bigl(\pi_{H,K}(\alpha)+\beta\bigr).
\]
Then the square
\begin{equation}\label{eq:affine-extension-universal-square}
\begin{tikzcd}[column sep=large]
F^-\times Y \arrow[r,"{(x,y)\mapsto x+d(y)}"] \arrow[d,"{\mathrm{pr}_Y}"']
  & F^+ \arrow[d,"{q_H}"]\\
Y \arrow[r,"{y\mapsto\alpha+q_Hd(y)}"']
  & \operatorname{Fib}_{\pi_{H,K}}\!\bigl(\pi_{H,K}(\alpha)+\beta\bigr)
\end{tikzcd}
\end{equation}
is a pullback.  Consequently its upper horizontal map is a regular epimorphism.  More generally, if $r:X\twoheadrightarrow F^-$ is any regular epimorphism, then
\[
X\times Y\longrightarrow F^+,
\qquad
(x,y)\longmapsto r(x)+d(y)
\]
is a regular epimorphism.
\end{lemma}

\begin{proof}
The lower arrow in \eqref{eq:affine-extension-universal-square} is the composite of the regular epimorphism \eqref{eq:affine-increment-regular-epi} with the isomorphism \eqref{eq:fibre-translation-isomorphism}.  The square is a pullback: for $y\in Y$ and $z\in F^+$ with
\[
q_H(z)=\alpha+q_Hd(y),
\]
the element $x=z-d(y)$ belongs to $F^-$ and is the unique point mapping to $(y,z)$.  Hence the upper map is the pullback of a regular epimorphism and is therefore a regular epimorphism.  The final assertion follows by composing it with the product regular epimorphism $r\times1_Y$.
\end{proof}

\begin{corollary}[Compatible affine extension]\label{cor:compatible-affine-extension}
Let $\mathscr A$ and $\mathscr B$ be finite parameter sets with label maps
\[
\lambda_{\mathscr A}:\mathscr A\to\mathcal C,
\qquad
\lambda_{\mathscr B}:\mathscr B\to\mathcal C,
\]
and suppose every pair of labels $\lambda_{\mathscr A}(s)$ and $\lambda_{\mathscr B}(b)$ is compatible.  Suppose
\[
\overline W\lambda_{\mathscr A}:\mathscr A\twoheadrightarrow\operatorname{Fib}_H(\alpha)
\]
is a regular epimorphism, $q_K\overline W\lambda_{\mathscr B}$ is constant with value $\beta\in A/K$, and
\begin{equation}\label{eq:compatible-affine-quotient-epi}
q_H\overline W\lambda_{\mathscr B}:\mathscr B\twoheadrightarrow
\operatorname{Fib}_{\pi_{H,K}}(\beta)
\end{equation}
is a regular epimorphism.  Then the compatible-union map
\[
\mathscr A\times\mathscr B\longrightarrow\mathcal C,
\qquad
(s,b)\longmapsto\lambda_{\mathscr A}(s)\cup\lambda_{\mathscr B}(b)
\]
has residue image exactly
\[
\operatorname{Fib}_K\!\bigl(\pi_{H,K}(\alpha)+\beta\bigr).
\]
If $\nu\lambda_{\mathscr B}$ is constant with value $c$, the component-count set is translated by $c$, so its spread is unchanged.
\end{corollary}

\begin{proof}
Apply Lemma~\ref{lem:affine-torsor-extension} to the regular epimorphism $\overline W\lambda_{\mathscr A}$ and the increment map $d=\overline W\lambda_{\mathscr B}$.  Compatibility makes the union map defined on the whole product, and \eqref{eq:observation-square} identifies its residue map with addition.
\end{proof}

\paragraph{Keeping the component count.}
Let $I,I'\subseteq\N$ be finite intervals, let $\alpha\in A/H$ and $\beta\in A/K$, and let $\mathscr L$ be a finite label-parameter set with label map $\lambda:\mathscr L\to\mathcal C$.  Suppose $q_K\overline W\lambda$ is constantly $\beta$.  Form
\[
Y_{I'}=\{(L,i)\in\mathscr L\times I:i+\nu(\lambda(L))\in I'\}
\]
and the response map
\begin{equation}\label{eq:graded-quotient-regular-epi}
\begin{aligned}
\Psi:Y_{I'}&\longrightarrow
\operatorname{Fib}_{\pi_{H,K}}\!\bigl(\pi_{H,K}(\alpha)+\beta\bigr)\times I',\\
(L,i)&\longmapsto
\bigl(\alpha+q_H\overline W(\lambda(L)),\,i+\nu(\lambda(L))\bigr).
\end{aligned}
\end{equation}
Assume $\Psi$ is a regular epimorphism and put
\[
T^-=\operatorname{Fib}_H(\alpha)\times I,
\qquad
T^+=\operatorname{Fib}_K\!\bigl(\pi_{H,K}(\alpha)+\beta\bigr)\times I'.
\]
Let
\[
\chi:T^+\longrightarrow
\operatorname{Fib}_{\pi_{H,K}}\!\bigl(\pi_{H,K}(\alpha)+\beta\bigr)\times I',
\qquad
(z,k')\longmapsto(q_H(z),k').
\]
Define the edge object by the pullback
\begin{equation}\label{eq:graded-local-translation}
\begin{tikzcd}[column sep=large]
E_{\operatorname{Tr}_{\mathscr L}} \arrow[r] \arrow[d,"{t}"']
  & Y_{I'} \arrow[d,"{\Psi}"]\\
T^+ \arrow[r,"{\chi}"']
  & \operatorname{Fib}_{\pi_{H,K}}\!\bigl(\pi_{H,K}(\alpha)+\beta\bigr)\times I'.
\end{tikzcd}
\end{equation}
For a point $((z,k'),(L,i))$ of this pullback put
\[
s((z,k'),(L,i))=(z-\overline W(\lambda(L)),i),
\qquad
\ell((z,k'),(L,i))=\lambda(L).
\]
\begin{corollary}[Graded affine translation]\label{cor:graded-affine-translation}
For the preceding pullback, $s$ lands in $T^-$, and $(s,\ell,t)$ is a covering labelled correspondence
\[
\operatorname{Tr}_{\mathscr L}(T^-,T^+):T^-\rightsquigarrow T^+.
\]
\end{corollary}

\begin{proof}
The pullback identity gives
\[
q_H(z)=\alpha+q_H\overline W(\lambda(L)),
\qquad
k'=i+\nu(\lambda(L)).
\]
Hence $z-\overline W(\lambda(L))\in\operatorname{Fib}_H(\alpha)$, so the source is well defined, and the conservation identity is immediate.  Since $\Psi$ is a regular epimorphism, pullback stability makes $t$ a regular epimorphism.
\end{proof}

\subsection{The prime-power filtration}

\begin{lemma}\label{lem:torsion-filtration}
For $m,n\ge1$,
\[
H_n=\langle[1/n]\rangle,
\qquad |H_n|=n,
\]
\[
H_m\le H_n\iff m\mid n,
\qquad
H_m\vee H_n=H_{\operatorname{lcm}(m,n)}.
\]
Consequently the nonzero join-irreducible stages are exactly $H_{p^e}$, with $p$ prime and $e\ge1$.
\end{lemma}

\begin{proof}
If $[q]\in H_n$, then $nq\in\Z$, so $[q]=[j/n]$ for some $j\in\Z$; hence $H_n=\langle[1/n]\rangle$ and has order $n$.  The inclusion criterion follows by testing $[1/m]$.  Put $L=\operatorname{lcm}(m,n)$.  Since $\gcd(L/m,L/n)=1$, there exist $a,b\in\Z$ with
\[
aL/m+bL/n=1.
\]
Dividing by $L$ gives $[1/L]\in H_m+H_n$, whence $H_m\vee H_n=H_L$.  The join-irreducibles of the divisibility lattice are precisely the prime powers.
\end{proof}

Accordingly we index the nontrivial steps by prime powers $Q=p^e$.  The \emph{rank} of the stage $H_{p^e}$ means the positive integer $Q=p^e$ itself.  Let $F_Q$ be the join of all prime-power stages of rank at most $Q$, and $F_{<Q}$ the join of those of smaller rank.  The simple quotient
\begin{equation}\label{eq:simple-quotient}
S_Q=F_Q/F_{<Q}
\end{equation}
has order $p$: the $p$-adic valuation of the lower least common multiple is $e-1$, and adjoining $p^e$ raises it by one.  Thus $S_Q$ is a one-dimensional $\F_p$-object.

For $X\ge1$ define the endpoint stage
\begin{equation}\label{eq:endpoint-stage}
E_X=H_{\operatorname{lcm}(1,\ldots,\lfloor X\rfloor)}.
\end{equation}
Every class $[a/b]\in A$ belongs to $H_b$, so
\begin{equation}\label{eq:cofinal-union}
A=\bigcup_{X\ge1}E_X.
\end{equation}
Equivalently, for every $B$ one has the identity of subsets of $A/E_B$
\begin{equation}\label{eq:centre-kernel-union}
A/E_B
=
\bigcup_{X\ge B}\ker\!\left(\pi_{B,X}:A/E_B\to A/E_X\right).
\end{equation}
Indeed $\ker(\pi_{B,X})=q_{E_B}(E_X)$, while \eqref{eq:cofinal-union} gives
\[
q_{E_B}(A)
=q_{E_B}\!\left(\bigcup_{X\ge B}E_X\right)
=\bigcup_{X\ge B}q_{E_B}(E_X).
\]
Consequently every affine centre is killed on a final segment of the endpoint filtration.

Finally, if $Q$ is a prime-power step, then
\begin{equation}\label{eq:endpoint-current-step}
E_{Q-1}=F_{<Q},
\qquad
E_Q=F_Q.
\end{equation}
Indeed a prime-power rank is $<Q$ exactly when it is at most $Q-1$.  Thus the local simple quotient \eqref{eq:simple-quotient} is literally the associated quotient of the global endpoint filtration at $Q$.
